\section{Programma}
\label{sec:Programma}

Di seguito viene riportato il codice sorgente del programma utilizzato per la trattazione.

\subsection{Contratto AbstractParam}
\label{subsec:AbstractParam}

Per il corretto utilizzo del programma e' necessario creare una classe che estende \textbf{AbstractParam}:
\begin{lstlisting}[style=myStile, caption={AbstractParam}, label={script:1}, escapeinside={(*@}{@*)}]
class AbstractParam(ABC):
    """
    Contratto obbligatorio per i Parametri
    """

    # ===== PARAMETRI BASE =====
    
    @property
    @final
    def n_roads(self) -> int:           #(*@$N$@*)
        """
        Numero totale di strade del grafo.
        
        Returns:
            int: numero di righe di Gamma_hat (una per ogni strada)
        """
        return self.Gamma_hat.shape[0]

    @property
    @final
    def n_routes_hat(self) -> int:      #(*@$\widehat{n}$@*)
        """
        Numero di percorsi disponibili per la popolazione hat.
        
        Returns:
            int: numero di colonne di Gamma_hat
        """
        return self.Gamma_hat.shape[1]

    @property
    @final
    def n_routes_check(self) -> int:    #(*@$\widecheck{n}$@*)
        """
        Numero di percorsi disponibili per la popolazione check.
        
        Returns:
            int: numero di colonne di Gamma_check
        """
        return self.Gamma_check.shape[1]

    @property
    @abstractmethod
    def operation(self) -> Operations:
        """
        Tipo di operazione da eseguire (es. Nash Equilibrium).
        
        Returns:
            Operations: tipo di operazione
        """
        pass

    @property
    @abstractmethod
    def output_directory(self) -> str:
        """
        Directory dove salvare i risultati.
        
        Returns:
            str: percorso della directory di output
        """
        pass

    @property
    def ymin_main_graph(self):
        """
        Valore minimo visualizzato nei grafici dei tempi di viaggio
        
        Returns:
            return[0]: valore minimo hat
            return[1]: valore minimo check
        """
        return [0., 0.]
    
    @property
    def ymax_main_graph(self):
        """
        Valore massimo visualizzato nei grafici dei tempi di viaggio
        
        Returns:
            return[0]: valore minimo hat
            return[1]: valore minimo check
        """
        return [0., 0.]

    @property
    def save_result(self) -> bool:
        """
        Indica se salvare i risultati su file.
        
        Returns:
            bool: True se i risultati devono essere salvati
        """
        return False
    
    @property
    def show_iterations(self) -> bool:
        """
        Indica se mostrare le iterazioni durante il calcolo.
        
        Returns:
            bool: True se si vogliono stampare a terminale le iterazioni
        """
        return True

    @property
    def save_as_fraction(self) -> bool:
        """
        Indica se salvare i valori come frazioni.
        Utilizzato solo nel caso in cui self.operation==Operations.NASH_EQ
        
        Returns:
            bool: True se i numeri devono essere convertiti in frazioni
        """
        return True

    @property
    def MIN(self) -> float:
        """
        Valore minimo della variazione nei grafici.
        
        Returns:
            float: limite inferiore
        """
        return 0.

    @property
    def MAX(self) -> float:
        """
        Valore massimo della variazione nei grafici.
        
        Returns:
            float: limite superiore
        """
        return 1.

    @property
    def step(self) -> float:
        """
        Passo della variazione nei grafici.
        
        Returns:
            float: incremento della variazione
        """
        return 0.01

    @property
    def entity_number_hat(self) -> float:   #(*@$\widehat{\zeta}$@*)
        """
        Numero totale di veicoli nella popolazione hat.
        Nel caso in cui self.operation==Operations.NASH_EQ_POP_VARIATIONS, viene
        utilizzato solo se self.variate_pop_hat==False
        
        Returns:
            float: numero di veicoli hat
        """
        return 1.
    
    @property
    def entity_number_check(self) -> float: #(*@$\widecheck{\zeta}$@*)
        """
        Numero totale di veicoli nella popolazione check. 
        Nel caso in cui self.operation==Operations.NASH_EQ_POP_VARIATIONS, viene
        utilizzato solo se self.variate_pop_check==False
                
        Returns:
            float: numero di veicoli check
        """
        return 1.

    @property
    def variation_coefficient_pop_hat(self) -> float:
        """
        Indica il coefficiente di variazione del numero di veicoli della popolazione hat durante il calcolo,
        nel caso in cui self.operation==Operation.NASH_EQ_POP_VARIATIONS
        
        Returns:
            float: coefficiente di variazione della popolazione hat
        """
        return 1.
    
    @property
    def variation_coefficient_pop_check(self) -> float:
        """
        Indica il coefficiente di variazione del numero di veicoli della popolazione check durante il calcolo,
        nel caso in cui self.operation==Operation.NASH_EQ_POP_VARIATIONS
        
        Returns:
            float: coefficiente di variazione della popolazione check
        """
        return 1.



    # ===== MATRICI =====
    
    @property
    @abstractmethod
    def Gamma_hat(self) -> SafeArray:   #(*@$\widehat{\Gamma}$@*)
        """
        Matrice dei percorsi per la popolazione hat.
        
        Returns:
            SafeArray: matrice (n_roads x n_routes_hat)
        """
        pass

    @property
    @abstractmethod
    def Gamma_check(self) -> SafeArray: #(*@$\widecheck{\Gamma}$@*)
        """
        Matrice dei percorsi per la popolazione check.
        
        Returns:
            SafeArray: matrice (n_roads x n_routes_check)
        """
        pass

    @property
    def variation_hat(self) -> SafeArray:
        """
            Ritorna array di dimensione (*@$N$@*), i cui elementi sono 1 o 0:
                - 'return[i]'!=0: se viene utilizzata una variazione, i costi della strada 'i' variano per la popolazione (*@$\textit{check}$@*) di un coefficiente 'return[i]', maggiore o minore di 0
                - 'return[i]'==0: se viene utilizzata una variazione, i costi della strada 'i' non variano per la popolazione (*@$\textit{hat}$@*)
        """
        return SafeArray(np.zeros(self.n_roads))

    @property
    def variation_check(self) -> SafeArray:
        """
            Ritorna array di dimensione (*@$N$@*), i cui elementi sono:
                - 'return[i]'!=0: se viene utilizzata una variazione, i costi della strada 'i' variano per la popolazione (*@$\textit{check}$@*) di un coefficiente 'return[i]', maggiore o minore di 0
                - 'return[i]'==0: se viene utilizzata una variazione, i costi della strada 'i' non variano per la popolazione (*@$\textit{check}$@*)
        """
        return SafeArray(np.zeros(self.n_roads))



    # ===== COSTI =====
    
    @abstractmethod
    def tau_hat(self, eta_hat, eta_check) -> SafeArray:     #(*@$\widehat{\tau}(\widehat{\eta}, \widecheck{\eta})$@*)
        """
        Calcola il costo (tempo di viaggio) per la popolazione hat.
        
        Args:
            eta_hat: veicoli hat presenti su ogni strada
            eta_check: veicoli check presenti su ogni strada
        
        Returns:
            SafeArray: costo per ogni strada (dimensione n_roads)
        """
        pass

    @abstractmethod
    def tau_check(self, eta_hat, eta_check) -> SafeArray:   #(*@$\widecheck{\tau}(\widehat{\eta}, \widecheck{\eta})$@*)
        """
        Calcola il costo (tempo di viaggio) per la popolazione check.
        
        Args:
            eta_hat: veicoli hat presenti su ogni strada
            eta_check: veicoli check presenti su ogni strada
        
        Returns:
            SafeArray: costo per ogni strada (dimensione n_roads)
        """
        pass
    
    @final
    def tau_hat_variated(self, eta_hat, eta_check, variation) -> SafeArray:
        """
        Calcola il costo per la popolazione hat includendo una variazione lineare.
        
        Il costo restituito è:
            tau_hat + variation_hat * variation
        
        Args:
            eta_hat: veicoli hat presenti su ogni strada
            eta_check: veicoli check presenti su ogni strada
            variation: coefficiente scalare di variazione
        
        Returns:
            SafeArray: costo modificato per ogni strada
        """
        return self.tau_hat(eta_hat, eta_check) + self.variation_hat * variation

    @final
    def tau_check_variated(self, eta_hat, eta_check, variation) -> SafeArray:
        """
        Calcola il costo per la popolazione check includendo una variazione lineare.
        
        Il costo restituito è:
            tau_check + variation_check * variation
        
        Args:
            eta_hat: veicoli hat presenti su ogni strada
            eta_check: veicoli check presenti su ogni strada
            variation: coefficiente scalare di variazione
        
        Returns:
            SafeArray: costo modificato per ogni strada
        """
        return self.tau_check(eta_hat, eta_check) + self.variation_check * variation
\end{lstlisting}
Le proprietà e i metodi nel \autoref{script:1} contrassegnati da \textbf{@abstractmethod} devono essere istanziati nella classe dei parametri creata dall'utente, quelli
contrassegnati da \textbf{@final} non devono essere instanziati, mentre i restanti hanno valori base che possono essere reistanziati o meno.

\subsection{Utility}
\label{subsec:Utility}

Viene introdotta una Enum per permettere all'utente di definire quale operazione eseguire (tramite \autoref{script:1}):
\begin{lstlisting}[style=myStile, caption={Operations}, label={script:2}]
  class Operations(Enum):
    #Calcolo dell'Equilibrio di Nash
    NASH_EQ = 1 

    #Calcolo dei grafici del tempo di viaggio e dell'Equilibrio di Nash dipendenti dalla variazione dei costi delle strade
    NASH_EQ_TIME_VARIATIONS = 2

    #Calcolo del grafico del tempo di viaggio e dell'Equilibrio di Nash dipendenti dalla variazione delle popolazioni
    NASH_EQ_POP_VARIATIONS = 3
\end{lstlisting}

Estendiamo anche la classe \textbf{np.ndarray} (vedi \cite{numpy_docs}) per gestire i casi in cui un valore di un array/matrice posto ad \textbf{np.inf} viene moltiplicato per 0.
Ciò può capitare soprattutto durante il calcolo dell'\autoref{eq:2}

\begin{lstlisting}[style=myStile, caption={SafeArray}, label={script:3}]
  import numpy as np

  class SafeArray(np.ndarray):
    def __new__(cls, input_array):
        obj = np.asarray(input_array, dtype=float).view(cls)
        return obj
    
    def __matmul__(self, other):
        selfChanged = False
        otherChanged = True        
        
        if np.isinf(self).any():
            self[np.isinf(self)] = 1.e200
            selfChanged = True
        if np.isinf(other).any():
            other[np.isinf(other)] = 1.e200
            otherChanged = False
            
        result = np.matmul(self, other)
        
        result[result > 1.e100] = np.inf
        
        if selfChanged:
            self[self>1.e100] = np.inf
        if otherChanged:
            other[other>1.e100] = np.inf

        return result.view(type(self))
\end{lstlisting}


\subsection{Classe Initializer}
\label{subsec:Initializer}

Una volta creata una classe che estende \textbf{AbstractParam}, per eseguire il programma è sufficiente creare un'istanza di \textbf{Initializer} e invocare il metodo 
\textbf{Initializer::start}:

\begin{lstlisting}[style=myStile, caption={Initializer}, label={script:4}, escapeinside={(*@}{@*)}]
class Initializer:
    
    def __init__(self, param: AbstractParam):
        np.seterr(divide='ignore')
        
        if not isinstance(param, AbstractParam):
            raise TypeError("param deve essere un'istanza di BaseParameters")
        
        self.param = param
        self.computer = Computer(self.param) #Vedi (*@\autoref{script:8}@*)
        self.initial_theta_hat = SafeArray(np.ones((self.param.n_routes_hat, 1))*(1/self.param.n_routes_hat))
        self.initial_theta_check = SafeArray(np.ones((self.param.n_routes_check, 1))*(1/self.param.n_routes_check))
        self.limit_hat = 1e-11
        self.limit_check = 1e-11
        
        
    def start(self):
        """
            In base a operation in param (vedi (*@\autoref{script:1}@*)), decide se calcolare l'Equilibrio di Nash (senza variazioni) o costruire il grafico
            variando i costi delle strade
        """
        if self.param.operation == Operations.NASH_EQ:
            return self.getNashEquilibria()
        else:
            return self.graphVariations()
\end{lstlisting}

Il metodo \textbf{Initializer::getNashEquilibria} calcola l'Equilibrio di Nash utilizzando i dati forniti tramite \textbf{param} (istanza di una classe figlia di \autoref{script:1}), senza variare
i costi delle strade e le veicoli nelle popolazioni. Decide inoltre se salvare il risultati su file verificando il valore attribuito a \textbf{param.save\_result}. Ritorna, in ogni
caso, i valori che descrivono l'Equilibrio di Nash della rete.

\begin{lstlisting}[style=myStile, caption={Initializer::getNashEquilibria}, label={script:5}, escapeinside={(*@}{@*)}]
def getNashEquilibria(self):
    """
        Calcola l'Equilibrio di Nash e ritorna la distribuzione della popolazione hat e check sulla rete, insieme ai tempi di viaggio della stessa.
        Se save_result==True, stampa tali risultati
    """     
    #Equilibrio di Nash:
    theta_hat, theta_check = self.computer.getNashEquilibria(self.initial_theta_hat, self.initial_theta_check, self.limit_hat, self.limit_check)
    
    #Tempi di viaggio su ogni strada:
    T_hat = self.computer.T_hat(theta_hat, theta_check, self.param.entity_number_hat, self.param.entity_number_check, 0)
    T_check = self.computer.T_check(theta_hat, theta_check, self.param.entity_number_hat, self.param.entity_number_check, 0)
    
    if self.param.save_result: self.saveNashEq(theta_hat, theta_check, T_hat, T_check)
    
    return theta_hat, theta_check, T_hat, T_check
\end{lstlisting}

Con: 

\begin{lstlisting}[style=myStile, caption={Initializer::saveNashEq}, label={script:14}, escapeinside={(*@}{@*)}]
def saveNashEq(self, theta_hat, theta_check, T_hat, T_check, file_name="nash_eq.txt"):
    content = self._format_nash_eq(theta_hat, theta_check, T_hat, T_check)

    file_path = Path(self.param.output_directory) / file_name

    with open(file_path, "w", encoding="utf-8") as f:
        f.write(content)

def _format_nash_eq(self, theta_hat, theta_check, T_hat, T_check):
    vec_frac = np.vectorize(
        lambda x: str(x) if np.isinf(x) or not self.param.save_as_fraction 
        else str(Fraction(x).limit_denominator(1000000))
    )

    lines = []
    lines.append("Equilibrio di Nash:")
    lines.append(f"theta_hat:\n{vec_frac(theta_hat)}")
    lines.append(f"theta_check:\n{vec_frac(theta_check)}")
    lines.append("")
    lines.append("Tempi di viaggio dei percorsi:")
    lines.append(f"T_hat:\n{vec_frac(T_hat)}")
    lines.append(f"T_check:\n{vec_frac(T_check)}")

    return "\n".join(lines)
\end{lstlisting}

Il metodo \textbf{Initializer::graphVariations} calcola i parametri utili alla costruzione dei grafici rappresentanti la variazione del tempo di 
viaggio e la variazione dell'Equilibrio di Nash rispetto alla variazione dei costi delle strade o del numero di veicoli delle 
popolazioni definite in \textbf{param}, \autoref{script:1}. Decide se costruire tali grafici e salvarli in \textbf{param.output\_directory}
verificando il valore attribuito a \textbf{param.save\_result}. Ritorna, in ogni caso, i valori utili a disegnarli.

\begin{lstlisting}[style=myStile, caption={Initializer::graphVariations}, label={script:6}]
    def graphVariations(self):
        """
        Calcola gli array per costruire i grafici del tempo di viaggio della rete e dell'Equilibrio di Nash
        per le due popolazioni in funzione della variazione (costi o popolazione).
        Restituisce gli array per costruire i grafici e opzionalmente salva i risultati.
        """
        p = self.param
        variation_values = np.arange(p.MIN, p.MAX + p.step, p.step)
        N = len(variation_values)

        time_hat = np.zeros(N)
        time_check = np.zeros(N)
        thetas_hat = np.zeros((p.n_routes_hat, N))
        thetas_check = np.zeros((p.n_routes_check, N))

        theta_hat, theta_check = self.initial_theta_hat, self.initial_theta_check

        for idx, v in enumerate(variation_values):
            if p.show_iterations:
                print(f"Current step: {idx+1}/{N}")

            #Vedi (*@\autoref{script:13}@*)
            theta_hat, theta_check, T_hat, T_check = self._compute_step(
                theta_hat, theta_check, v
            )

            i_hat = np.argmax(theta_hat > 0)
            i_check = np.argmax(theta_check > 0)

            time_hat[idx] = T_hat[i_hat]
            time_check[idx] = T_check[i_check]
            thetas_hat[:, idx] = theta_hat[:, 0]
            thetas_check[:, idx] = theta_check[:, 0]

        if p.save_result:
            label = ('Variazione del costo delle strade'
                    if p.operation == Operations.NASH_EQ_TIME_VARIATIONS
                    else 'Variazione del numero di individui nella popolazione')

            #Grafico del tempo di viaggio della rete in funzione dalla variazione
            self.saveResultVariation(variation_values, time_hat, time_check, label)
            
            #Grafici dell'Equilibrio di Nash in funzione della variazione
            self.saveResultThetaVariation(variation_values, thetas_hat, 'hat', label)
            self.saveResultThetaVariation(variation_values, thetas_check, 'check', label)

        return variation_values, time_hat, time_check, thetas_hat, thetas_check
\end{lstlisting}

Con:
\begin{lstlisting}[style=myStile, caption={Initializer::\_compute\_step}, label={script:13}, escapeinside={(*@}{@*)}]
def _compute_step(self, theta_hat, theta_check, v):
    p = self.param
    c = self.computer

    if p.operation == Operations.NASH_EQ_TIME_VARIATIONS:
        theta_hat, theta_check = c.getNashEquilibriaCostVariation(  #Vedi (*@\autoref{script:10}@*)
            theta_hat, theta_check, self.limit_hat, self.limit_check, v
        )
        pop_hat, pop_check, var = p.entity_number_hat, p.entity_number_check, v

    else:
        pop_hat = p.entity_number_hat + (v*p.variation_coefficient_pop_hat)
        pop_check = p.entity_number_check + (v*p.variation_coefficient_pop_check)

        theta_hat, theta_check = c.getNashEquilibriaPopVariation(   #Vedi (*@\autoref{script:10}@*)
            theta_hat, theta_check, self.limit_hat, self.limit_check, pop_hat, pop_check
        )
        var = 0

    T_hat = c.T_hat(theta_hat, theta_check, pop_hat, pop_check, var)
    T_check = c.T_check(theta_hat, theta_check, pop_hat, pop_check, var)

    return theta_hat, theta_check, T_hat, T_check
\end{lstlisting}

Per disegnare i grafici calcolati nel \autoref{script:6} si fa uso della libreria \textbf{matplotlib.pyplot} (vedi \cite{matplotlib_docs}).

\begin{lstlisting}[style=myStile, caption={Initializer::showResultPlot}, label={script:7}]
    import matplotlib.pyplot as plt

    ...

    def saveResultVariation(self, variation_values, time_of_travel_hat, time_of_travel_check, xLabel):        
        # Grafico Hat
        plt.figure()
        plt.plot(variation_values, time_of_travel_hat, label='Hat', color='blue')
        plt.xlabel(xLabel)
        plt.ylabel('Tempo di viaggio')
        plt.title(f'Tempo di viaggio (Hat)\nin funzione di: {xLabel}')
        plt.legend()
        plt.grid(True)
        
        if self.param.ymin_main_graph[0] != self.param.ymax_main_graph[0]:
            plt.ylim(self.param.ymin_main_graph[0], self.param.ymax_main_graph[0])
        
        self.saveGraph(f'main_hat_{xLabel}')
        
        

        # Grafico Check
        plt.figure()
        plt.plot(variation_values, time_of_travel_check, label='Check', color='red')
        plt.xlabel(xLabel)
        plt.ylabel('Tempo di viaggio')
        plt.title(f'Tempo di viaggio (Check)\nin funzione di: {xLabel}')
        plt.legend()
        plt.grid(True)
        
        if self.param.ymin_main_graph[1] != self.param.ymax_main_graph[1]:
            plt.ylim(self.param.ymin_main_graph[1], self.param.ymax_main_graph[1])
            
        self.saveGraph(f'main_check_{xLabel}')
        
    def saveResultThetaVariation(self, variation_values, thetas, pop, xLabel):
        n = thetas.shape[0]

        subfolder = "hat" if pop == "hat" else "check"
        color = "blue" if pop == "hat" else "red"

        for i in range(n):
            plt.figure(figsize=(6, 4))

            plt.plot(
                variation_values,
                thetas[i, :],
                linewidth=1.8,
                color=color
            )

            plt.title(f'Theta_{pop}_{i+1}')
            plt.xlabel(xLabel)
            plt.ylabel('Distribuzione popolazione')
            plt.ylim(0, 1.2)
            plt.grid(True)

            file_label = f"{pop}/{subfolder}_theta_{i+1}_{xLabel}"

            self.saveGraph(file_label)

            plt.close()
 
 
    def saveGraph(self, xLabel):
        file_name = f"grafico_{xLabel}.png".replace(" ", "_")

        file_path = Path(self.param.output_directory) / file_name
        file_path.parent.mkdir(parents=True, exist_ok=True)

        plt.savefig(file_path, dpi=300, bbox_inches='tight', pad_inches=0.5)
\end{lstlisting}

\subsection{Classe Computer}
\label{subsec:Computer}

Per calcolare gli Equilibri di Nash e' necessario creare un'istanza della classe \textbf{Computer}. Tale operazione non deve, durante un utilizzo standard, essere effettuata dall'utente,
bensì viene mascherata ad esso dall'utilizzo di \textbf{Initializer}, \autoref{script:4}

\begin{lstlisting}[style=myStile, caption={Computer}, label={script:8}]
    class Computer:    
    def __init__(self, param: AbstractParam):
        self.param = param
        
        #Parametro utile durante il calcolo di f_hat e f_check, vedi (*@\autoref{script:12}@*)
        self.lmbda = 0.5*np.min([1/self.param.n_routes_hat, 1/self.param.n_routes_check])
        
        #Array colonna di 1 con dimensioni n_routes_hat ed n_routes_check 
        #Utile durante il calcolo di f_hat e f_check, vedi (*@\autoref{script:12}@*)
        self.one_n_hat = SafeArray(np.ones((self.param.n_routes_hat, 1)))
        self.one_n_check = SafeArray(np.ones((self.param.n_routes_check, 1)))
\end{lstlisting}

Per calcolare l'Equilibrio di Nash sfruttiamo la dimostrazione (vedi \cite[Cap. 4]{MR4911797}) del \autoref{th:1}: individuiamo un punto fisso della funzione $\mathcal{F}$.
Cominciamo da un $(\widehat{\theta}, \widecheck{\theta})\in S^{\widehat{n}}\times S^{\widecheck{n}}$ generico; si ricava $\mathcal{F}(\widehat{\theta}, \widecheck{\theta})$,
che va a sostituire il valore di $(\widehat{\theta}, \widecheck{\theta})$, salvato invece come $(\widehat{\theta}_{prev}, \widecheck{\theta}_{prev})$. Questa operazione viene ripetuta finché
$\norm{\widehat{\theta} - \widehat{\theta}_{prev}} < \widehat{limit}$ e $\norm{\widecheck{\theta} - \widecheck{\theta}_{prev}} < \widecheck{limit}$, dove i $limit$ sono limiti di precisione
definiti in \textbf{Initializer} (vedi \autoref{script:4}). Introduciamo anche un coefficiente di variazione e il numero di veicoli nelle popolazioni, utili per il calcolo effettuato
nel \autoref{script:13}, che verrà applicato ai costi delle strade durante il calcolo dei tempi di viaggio della rete, vedi \autoref{script:11}. 

\begin{lstlisting}[style=myStile, caption={Computer::getNashEquilibriaVariation}, label={script:9}]
def getNashEquilibriaVariation(self, theta_hat, theta_check, limit_hat, limit_check, zeta_hat, zeta_check, variationTime):
    """
        Ritorna l'Equilibrio di Nash variando i costi delle strade selezionate in param di un coefficiente variation, usando come
        punto di partenza theta_hat, theta_check, come limite di precisione limit_hat, limit_check e come numero di veicoli zeta_hat e zeta_check
        :param theta_hat, theta_check: vettore di dimensione [param.n_routes_hat, 1] e [param.n_routes_check, 1]
        :param limit_hat, limit_check, variation, zeta_hat, zeta_check: float
    """
    prev_theta_hat = SafeArray(np.ones((self.param.n_routes_hat, 1))*np.inf)
    prev_theta_check = SafeArray(np.ones((self.param.n_routes_check, 1))*np.inf)
    
    count = 0
    
    while np.linalg.norm(theta_hat-prev_theta_hat) > limit_hat or np.linalg.norm(theta_check-prev_theta_check) > limit_check:
        prev_theta_hat[:] = theta_hat
        prev_theta_check[:] = theta_check
        
        theta_hat = self.f_hat(theta_hat, theta_check, zeta_hat, zeta_check, variationTime)
        theta_check = self.f_check(theta_hat, theta_check, zeta_hat, zeta_check, variationTime)
        count+=1
    
    if self.param.show_iterations: print("Iterations: " + str(count))
    return theta_hat, theta_check
\end{lstlisting}

Introduciamo 2 metodi per calcolare l'Equilibrio di Nash in 2 casi specifici.
Durante il calcolo della variazione dell'Equilibrio di Nash in funzione della variazione dei costi delle strade viene utilizzato (vedi \autoref{script:13}):
\begin{lstlisting}[style=myStile, caption={Computer::getNashEquilibriaCostVariation}, label={script:15}]
    def getNashEquilibriaCostVariation(self, theta_hat, theta_check, limit_hat, limit_check, variation):
        """
            Ritorna l'Equilibrio di Nash variando i costi delle strade selezionate in param di un coefficiente variation, usando come
            punto di partenza theta_hat, theta_check, come limite di precisione limit_hat, limit_check e come numero di veicoli i self.param.entity_number
            :param theta_hat, theta_check: vettore di dimensione [param.n_routes_hat, 1] e [param.n_routes_check, 1]
            :param limit_hat, limit_check, variation, totPopulation: float
        """
        return self.getNashEquilibriaVariation(theta_hat, 
                                               theta_check, 
                                               limit_hat, 
                                               limit_check, 
                                               self.param.entity_number_hat, 
                                               self.param.entity_number_check, 
                                               variation)
\end{lstlisting}

Durante il calcolo della variazione dell'Equilibrio di Nash in funzione della variazione del numero di veicoli viene utilizzato (vedi \autoref{script:13}):
\begin{lstlisting}[style=myStile, caption={Computer::getNashEquilibriaCostVariation}, label={script:16}]
    def getNashEquilibriaPopVariation(self, theta_hat, theta_check, limit_hat, limit_check, zeta_hat, zeta_check):
        """
            Ritorna l'Equilibrio di Nash usando come costi delle strade i tau delle rispettive popolazioni presenti in param (senza variazioni), come
            punto di partenza theta_hat, theta_check, come limite di precisione limit_hat, limit_check e come numero di veicoli zeta_hat e zeta_check
            :param theta_hat, theta_check: vettore di dimensione [param.n_routes_hat, 1] e [param.n_routes_check, 1]
            :param limit_hat, limit_check, totPopulation, zeta_hat, zeta_check: float
        """
        return self.getNashEquilibriaVariation(theta_hat,
                                               theta_check,
                                               limit_hat,
                                               limit_check,
                                               zeta_hat,
                                               zeta_check,
                                               0)
\end{lstlisting}

Si può sfruttare la funzione \textbf{getNashEquilibria} per calcolare l'Equilibrio di Nash senza introdurre una variazione ai costi delle strade:

\begin{lstlisting}[style=myStile, caption={Computer::getNashEquilibria}, label={script:10}]
    def getNashEquilibria(self, theta_hat, theta_check, limit_hat, limit_check):
        """
            Ritorna l'Equilibrio di Nash usando come costi delle strade i tau delle rispettive popolazioni presenti in param (senza variazioni e con numero di veicoli
            nelle popolazioni 1), come punto di partenza theta_hat, theta_check e come limite di precisione limit_hat, limit_check
            :param theta_hat, theta_check: vettore di dimensione [param.n_routes_hat, 1] e [param.n_routes_check, 1]
            :param limit_hat, limit_check: float
        """
        return self.getNashEquilibriaVariation(theta_hat, 
                                               theta_check, 
                                               limit_hat, 
                                               limit_check, 
                                               self.param.entity_number_hat, 
                                               self.param.entity_number_check,
                                               0)
\end{lstlisting}

Vengono introdotte le seguenti funzioni ausiliarie:

\begin{lstlisting}[style=myStile, caption={Funzioni ausiliarie di Computer}, label={script:11}, escapeinside={(*@}{@*)}]
    #(*@\autoref{eq:4}@*)
    def phi(self, x):
        x = np.asarray(x)
        res = np.empty_like(x)

        mask = np.isinf(x)
        res[mask] = 1.
        res[~mask] = x[~mask] / (1. - x[~mask])

        return SafeArray(res.reshape(-1, 1))
    
    #(*@\autoref{eq:2}@*) utilizzando i tau come nell'(*@\autoref{eq:10}@*)
    def T_hat(self, theta_hat, theta_check, zeta_hat, zeta_check, variation):
        nu_hat = self.param.Gamma_hat @ theta_hat

        nu_check = self.param.Gamma_check @ theta_check

        return self.param.Gamma_hat.T @ self.param.tau_hat_variated(nu_hat*zeta_hat, nu_check*zeta_check, variation)
    
    #(*@\autoref{eq:2}@*) utilizzando i tau come nell'(*@\autoref{eq:10}@*)
    def T_check(self, theta_hat, theta_check, zeta_hat, zeta_check, variation):
        nu_hat = self.param.Gamma_hat @ theta_hat

        nu_check = self.param.Gamma_check @ theta_check

        return self.param.Gamma_check.T @ self.param.tau_check_variated(nu_hat*zeta_hat, nu_check*zeta_check, variation)
    
    #(*@\autoref{eq:5}@*)
    def normalize(self, theta):
        theta_max = np.maximum(theta, 0)
        sum = np.sum(theta_max)
        return theta_max / sum
\end{lstlisting}

Infine, vengono definite le funzioni per il calcolo di $\widehat{\mathcal{F}}$ e $widecheck{\mathcal{F}}$ (\autoref{eq:7}):

\begin{lstlisting}[style=myStile, caption={Computer::f\_hat, Computer::f\_check}, label={script:12}]
    def f_hat(self, theta_hat, theta_check, zeta_hat, zeta_check, variation):
        composition = self.phi(self.T_hat(theta_hat, theta_check, zeta_hat, zeta_check, variation))
        return self.normalize(theta_hat - self.lmbda*(composition - (theta_hat.T @ composition)[0]*self.one_n_hat))
    
    
    def f_check(self, theta_hat, theta_check, zeta_hat, zeta_check, variation):
        composition = self.phi(self.T_check(theta_hat, theta_check, zeta_hat, zeta_check, variation))
        return self.normalize(theta_check - self.lmbda*(composition - (theta_check.T @ composition)[0]*self.one_n_check))
\end{lstlisting}